% --- 霍尔效应实验报告正文 ---

\section{实验目的：}
\begin{enumerate}[label=\arabic* 、, leftmargin=4em, labelsep=0pt]
    \item 了解产生霍尔效应的物理过程及其测量磁场的原理、方法和实验误差产生的原因；
    \item 验证霍尔控制电流$I_s$与霍尔电压$U_H$的线性关系；
    \item 测定励磁电流$I_M$与霍尔电压$U_H$线性关系的范围。
\end{enumerate}

\section{实验原理：}

霍尔效应在1879年由美国物理学家埃德温·赫伯特·霍尔（Edwin Herbert Hall）发现。霍尔效应（Hall effect）是指当固体导体放置在一个磁场内，且有电流通过时，导体内的电荷载流子受到洛伦兹力而偏向一边，继而产生电压（霍尔电压）的现象。电压所引致的电场力会平衡洛伦兹力。通过霍尔效应，可证实导体内部的电流是由带有负电荷的粒子（自由电子）运动所造成的。霍尔效应广泛应用于传感器技术、电力电子、自动化控制和材料科学等多个领域，对测量和控制电流、磁场和电压具有重要作用。

霍尔效应是磁电效应的一种。在匀强磁场中放一块半导体薄板，使板面与磁场方向垂直，在半导体薄板中沿着与磁场垂直的方向通电流时，薄板的两个侧面间会出现电势差。半导体霍尔元件在磁测量中应用广泛，现在通用的特斯拉计，其探头就是霍尔元件。霍尔效应是研究半导体材料性能的重要理论依据，利用半导体材料制成的霍尔元件，已被广泛应用于科学技术领域和现代检测技术中，如磁场、位置、位移、转速等的测量。

\subsection{霍尔效应}

霍尔效应从本质上讲是运动的带电粒子在磁场中受洛伦兹力作用而引起的偏转。将半导体置于图5-1所示的磁场中。半导体长为$l$，宽为$b$，厚为$d$，外磁场$\mathbf{B}$沿$z$轴方向。若半导体中通有电流$I_s$，沿$x$轴方向，垂直于$\mathbf{B}$，则在垂直于$\mathbf{B}$和$I_s$方向，即沿$y$轴方向，出现一个横向电位差$U_H$，这一现象称霍尔效应，其横向电压$U_H$称霍尔电压。

% 此处可插入图5-1：霍尔效应示意图
% \begin{figure}[htbp]
%   \centering
%   \includegraphics[width=0.5\textwidth]{assets/fig5-1.png}
%   \caption{图5-1 霍尔效应}
% \end{figure}

设有横向均匀磁场$\mathbf{B}$作用在一束以速度为$\mathbf{v}$的电子流上，运动电荷受到的洛伦兹力为：
\begin{equation}
    \mathbf{F}_L = -e\mathbf{v} \times \mathbf{B} \label{eq:5-1}
\end{equation}
其方向沿$y$轴的负方向，造成电子流发生偏转，被推向半导体的上侧（$N$面）。在$N$面上形成负电荷积累，而相对下侧（$M$面）形成正电荷积累，使$MN$面之间建立了电场，其方向沿$y$轴正向，两侧的霍尔电压$U_H = E_H b$。

\subsection{霍尔电压$U_H$与外磁场$B$的关系}

电子在霍尔电场中同样受到一个静电场力：
\begin{equation}
    \mathbf{F}_E = -e\mathbf{E}_H
\end{equation}
沿$y$轴负向，它将阻碍电子向$M$面的积聚。随着积聚电荷的增加，电场不断增强，直到$F_E = F_L$，达到平衡，这时有：
\begin{equation}
    e\frac{U_H}{b} = evB
\end{equation}
设半导体薄片中电子浓度为$n$，则有：
\begin{equation}
    I_s = nevbd
\end{equation}
所以霍尔电压为：
\begin{equation}
    U_H = \frac{1}{ne}\frac{I_s B}{d} = K_H I_s B \label{eq:5-4}
\end{equation}
式中$R_H = \frac{1}{ne}$称为霍尔系数，它决定于材料的性质（种类、载流子的浓度），是反映材料霍尔效应强弱的重要参数。$K_H = \frac{R_H}{d} = \frac{1}{ned}$称为霍尔元件的灵敏度，它决定于材料的性质和几何尺寸，对于确定的霍尔片来说是一个常数，常采用单位 mV/(mA·T)。对于测磁场而言，$K_H$越大越好，因此常用半导体材料制成测磁场的霍尔传感器。$K_H$与霍尔片厚度$d$呈反比，所以霍尔元件都做得很薄，一般只有约0.26mm厚。

对于选定的霍尔片，保持通过它的工作电流$I_s$不变时，霍尔电压$U_H$与被测量的磁感应强度$B$呈正比，即：
\begin{equation}
    U_H = K_H I_s B \label{eq:5-5}
\end{equation}

若将测得的$U_H$值进行放大，用电表来表示，并通过一定的换算，在电表的表盘上直接以$B$值来标度，这就是霍尔效应测磁场的原理，由此原理制作的磁场测试仪也叫特斯拉计。

由于霍尔效应建立电场所需时间很短（约$10^{-12}\sim10^{-14}$s），因此通过霍尔元件的电流用直流或交流都可以。若霍尔元件的电流为交流：$I_s = I_0 \sin\omega t$，则所得的霍尔电压也是交流的：$U_H = K_H B I_0 \sin\omega t$。

\subsection{实验中的副效应及其消除方法}

事实上，在使用霍尔片测量时，得到的霍尔电压还包含了由不等势电压、温差电效应、能斯特效应、里吉-勒迪克（Righi-Leduc）效应等因素带来的附加电压，这里一一做简单介绍。

\textbf{(1) 不等势电压 $U_0$}

如图5-2所示，经常出现器件的1、2两电极的位置不在一个理想的等势面上的情况，因此即使不加磁场，只要有电流$I_s$通过，就有电压$U_0 = I_s r$产生，$r$为1、2两电极所在的两等势面之间的电阻。结果在测量$U_H$时，就叠加了$U_0$，使得$U_H$值偏大（当$U_0$与$U_H$同号）或偏小（当$U_0$与$U_H$异号）。显然$U_H$的符号取决于$I_s$和$B$两者的方向，而$U_0$只与$I_s$的方向有关，因此可以通过改变$I_s$的方向予以消除。

% 此处可插入图5-2：不等势电势差U₀

\textbf{(2) 温差电效应引起的附加电压 $U_E$}

如图5-3所示，由于构成电流的载流子速度不同，若速度为$v$的载流子所受的洛伦兹力与霍尔电场的作用力刚好抵消，则速度大于或小于$v$的载流子在电场和磁场作用下，将各自朝对立面偏转，从而在$y$方向引起温差$T_1-T_2$，由此产生的温差电效应。在1、2两电极上引入附加电压$U_E$，且其符号与$I_s$和$\mathbf{B}$的方向关系跟$U_H$是相同的，因此不能用改变$I_s$和$\mathbf{B}$方向的方法予以消除。但其引入的误差很小，可以忽略。

% 此处可插入图5-3：温差电效应

\textbf{(3) 能斯特（Nernst）效应引起的附加电压 $U_N$}

如图5-4所示，因器件两端电流引线的接触电阻不等，通电后在接点两处将产生不同的焦耳热，导致在$x$方向有温度梯度，引起载流子沿梯度方向扩散而产生热扩散电流。热流$Q$在$z$方向磁场作用下，类似于霍尔效应在$y$方向产生一附加电场$E_N$，相应的电压$U_N \propto I_M$，而$U_N$的符号只与$\mathbf{B}$的方向有关，与$I_s$的方向无关，因此可通过改变$\mathbf{B}$的方向予以消除。

% 此处可插入图5-4：能斯特效应

\textbf{(4) 里吉-勒迪克（Righi-Leduc）效应引起的附加电压 $U_{RL}$}

如图5-5所示，由于存在温度梯度$T_1'-T_2'$，由此引入的附加电压$U_{RL} \propto I_M$，$U_{RL}$的符号只与$\mathbf{B}$的方向有关，通过改变$\mathbf{B}$的方向，也能消除。

% 此处可插入图5-5：里吉-勒迪克效应

综上所述，实验中测得的1、2两电极的电压中，除$U_H$外还包含$U_0$、$U_N$、$U_{RL}$和$U_E$各电压的代数和，其中$U_0$、$U_N$和$U_{RL}$均可通过$I_s$和$B$的换向，即对称测量法（异号法）予以消除。四种工作状态测量的情况表示如下：

\begin{center}
\begin{tabular}{c|c}
\hline
\textbf{工作状态} & \textbf{测得电压} \\
\hline
$+I_s,\;+B$ & $U_1 = +U_H + U_0 + U_E + U_N + U_{RL}$ \\
$+I_s,\;-B$ & $U_2 = -U_H + U_0 - U_E - U_N - U_{RL}$ \\
$-I_s,\;-B$ & $U_3 = +U_H - U_0 + U_E - U_N - U_{RL}$ \\
$-I_s,\;+B$ & $U_4 = -U_H - U_0 - U_E + U_N + U_{RL}$ \\
\hline
\end{tabular}
\end{center}

对上述四式做如下运算，有：
\begin{equation}
    U_H = \frac{U_1 - U_2 + U_3 - U_4}{4} \label{eq:5-6}
\end{equation}

可见，除了温差电效应$U_E$以外的其他副效应产生的电压会全部消除。但是在非大电流、非强磁场下，$U_H \gg U_E$，因而$U_E$可以忽略不计。

一般情况下当$U_H$较大时，$U_1$与$U_3$同号，$U_2$与$U_4$同号，而两组数据反号，因此上式也可以写为：
\begin{equation}
    U_H = \frac{1}{4}\left(|U_1| + |U_2| + |U_3| + |U_4|\right) \label{eq:5-7}
\end{equation}
即用四次测量值的绝对值的平均值即可。

通过四种工作状态的换测消除伴随效应的方法，是消除系统误差的一种常用方法。

\section{实验仪器：}

实验室提供的霍尔效应实验组合仪由实验仪HL-A和测试仪HL-CF两部分组成。

\subsection{霍尔元件}

是由N型硅单晶经过硅平面工艺制成的磁电转换元件，元件尺寸大致为$0.3\mathrm{mm} \times 0.3\mathrm{mm} \times 0.26\mathrm{mm}$，元件胶合在绝缘衬板上，有四条引出线，其中两条导线为工作电流极，两条导线为霍尔电压输出极。同时将这4条引线焊接在绝缘衬板上，然后引到仪器换向开关上，能方便进行实验。霍尔元件的灵敏度一般在$10.0\mathrm{V/A}\cdot\mathrm{T}$左右，温度变化时，灵敏度也略有变化，这主要是由于不同温度下半导体的载流子浓度不同所造成。

\subsection{电磁铁}

采用冷轧电工钢制成铁芯，线圈用高强度漆包线多层密绕，层间绝缘，导线的绕向即激励电流的方向已标明在线圈上，可确定磁场方向。

\subsection{换向开关}

仪器装有三个换向开关，很方便地改变$I_s$、$I_M$、$U_H$的方向。

\subsection{测试仪面板说明}

HL-CF测试仪面板包含以下部分：
\begin{enumerate}[label=\arabic* 、, leftmargin=4em, labelsep=0pt]
    \item \textbf{工作电流输出}：为霍尔元件提供工作电流$I_s$，配有输出调节旋钮；
    \item \textbf{霍尔电压输入}：连接霍尔元件的电压输出端，测量霍尔电压$U_H$；
    \item \textbf{励磁恒流输出}：为电磁铁提供励磁电流$I_M$，配有输出调节旋钮。
\end{enumerate}

% 此处可插入图5-6：实验仪接线图
% \begin{figure}[htbp]
%   \centering
%   \includegraphics[width=0.6\textwidth]{assets/fig5-6.png}
%   \caption{图5-6 实验仪接线图}
% \end{figure}

\section{实验内容：}
\begin{enumerate}[label=\arabic* 、, leftmargin=4em, labelsep=0pt]
    \item 按图5-6连接实验仪器，将霍尔元件放置于电磁铁气隙中；
    \item 打开电源，调节工作电流$I_s$至适当值（如$I_s=3.00\text{mA}$）；
    \item 保持$I_s$不变，改变励磁电流$I_M$，分别测量四种工作状态（$+I_s,+B$；$+I_s,-B$；$-I_s,-B$；$-I_s,+B$）下的霍尔电压$U_1$、$U_2$、$U_3$、$U_4$；
    \item 按公式\eqref{eq:5-6}计算霍尔电压$U_H$，绘制$U_H-I_M$关系曲线，验证线性关系；
    \item 保持励磁电流$I_M$不变，改变工作电流$I_s$，测量对应的霍尔电压，验证$U_H$与$I_s$的线性关系。
\end{enumerate}

\section{注意事项：}
\begin{enumerate}[label=\arabic* 、, leftmargin=4em, labelsep=0pt]
    \item 实验前应仔细检查接线是否正确，特别注意工作电流和励磁电流的极性；
    \item 霍尔元件是易损器件，工作电流不能超过规定值（一般$I_s \leq 10\text{mA}$），否则会烧毁元件；
    \item 电磁铁的励磁电流不宜过大，以免线圈过热；
    \item 测量时应等待读数稳定后再记录数据；
    \item 利用对称测量法消除副效应时，必须严格按四种工作状态分别测量。
\end{enumerate}

\section{实验数据：}

\noindent 实验所用霍尔元件编号：No. 12-027。电磁铁线圈常数 $k = 261.5\;\text{G/A}$（即 $B = k \cdot I_M$）。

\subsection{保持 $I_s$ 不变，改变励磁电流 $I_M$}

固定工作电流 $I_s = 3.00\;\text{mA}$，改变励磁电流 $I_M$，采用对称测量法分别测量四种工作状态下的电压值 $U_1$、$U_2$、$U_3$、$U_4$，
按式\eqref{eq:5-6}计算霍尔电压 $U_H$。实验数据如表~\ref{tab:data1} 所示。

\begin{table}[htbp]
    \centering
    \caption{$I_s = 3.00\;\text{mA}$ 时，不同励磁电流下的霍尔电压测量值}
    \label{tab:data1}
    \begin{tabular}{cccccc}
        \hline
        $I_M$ (mA) & $U_1$ (mV) & $U_2$ (mV) & $U_3$ (mV) & $U_4$ (mV) & $U_H$ (mV) \\
        \hline
        200 &  13.9  & $-13.7$ &  13.7  & $-13.7$ & 13.75 \\
        400 &  27.7  & $-27.4$ &  27.4  & $-27.5$ & 27.50 \\
        600 &  41.3  & $-41.2$ &  41.2  & $-41.3$ & 41.25 \\
        800 &  55.1  & $-54.9$ &  54.9  & $-55.1$ & 55.00 \\
        1000&  68.8  & $-68.7$ &  68.7  & $-68.8$ & 68.75 \\
        \hline
    \end{tabular}
\end{table}

\noindent 其中 $U_H$ 按对称测量法公式计算，以 $I_M = 200\;\text{mA}$ 为例：
\begin{equation}
    U_H = \frac{U_1 - U_2 + U_3 - U_4}{4} = \frac{13.9 - (-13.7) + 13.7 - (-13.7)}{4} = \frac{55.0}{4} = 13.75\;\text{mV}
\end{equation}

\subsection{保持 $I_M$ 不变，改变工作电流 $I_s$}

固定励磁电流 $I_M = 600\;\text{mA}$，改变工作电流 $I_s$，同样采用对称测量法。实验数据如表~\ref{tab:data2} 所示。

\begin{table}[htbp]
    \centering
    \caption{$I_M = 600\;\text{mA}$ 时，不同工作电流下的霍尔电压测量值}
    \label{tab:data2}
    \begin{tabular}{cccccc}
        \hline
        $I_s$ (mA) & $U_1$ (mV) & $U_2$ (mV) & $U_3$ (mV) & $U_4$ (mV) & $U_H$ (mV) \\
        \hline
        1.00 &  13.9  & $-13.7$ &  13.7  & $-13.7$ & 13.75 \\
        3.00 &  41.3  & $-41.2$ &  41.2  & $-41.3$ & 41.25 \\
        5.00 &  68.8  & $-68.7$ &  68.7  & $-68.8$ & 68.75 \\
        7.00 &  96.3  & $-96.2$ &  96.2  & $-96.3$ & 96.25 \\
        9.00 & 123.8  & $-123.7$& 123.7  & $-123.8$& 123.75 \\
        \hline
    \end{tabular}
\end{table}

\section{数据分析：}

\subsection{$U_H$--$I_M$ 关系分析}

根据表~\ref{tab:data1} 数据，以励磁电流 $I_M$ 为横坐标、霍尔电压 $U_H$ 为纵坐标作图。
数据点 $(200, 13.75)$、$(400, 27.50)$、$(600, 41.25)$、$(800, 55.00)$、$(1000, 68.75)$ 呈现良好的线性关系。

对 $U_H$--$I_M$ 数据进行线性拟合（过原点），可得斜率：
\begin{equation}
    k_1 = \frac{\Delta U_H}{\Delta I_M} = \frac{68.75 - 13.75}{1000 - 200} = \frac{55.0}{800} = 0.06875\;\text{mV/mA}
\end{equation}

由式\eqref{eq:5-4} $U_H = K_H I_s B$ 及 $B = k \cdot I_M$（$k = 261.5\;\text{G/A} = 2.615 \times 10^{-2}\;\text{T/A}$），有：
\begin{equation}
    U_H = K_H I_s \cdot k I_M
\end{equation}

因此斜率 $k_1 = K_H I_s k$，可求得霍尔灵敏度：
\begin{equation}
    K_H = \frac{k_1}{I_s \cdot k} = \frac{0.06875 \times 10^{-3}\;\text{V/mA}}{3.00\;\text{mA} \times 2.615 \times 10^{-2}\;\text{T/A}} = \frac{6.875 \times 10^{-5}}{7.845 \times 10^{-2}} \approx 8.76 \times 10^{-4}\;\text{V/(mA}\cdot\text{T)}
\end{equation}

换算为常用单位：
\begin{equation}
    K_H \approx 0.876\;\text{mV/(mA}\cdot\text{T)} = 8.76 \times 10^{2}\;\text{mV/(A}\cdot\text{T)}
\end{equation}

\noindent 该灵敏度值在半导体霍尔元件的典型范围内，表明测量结果合理。

\subsection{$U_H$--$I_s$ 关系分析}

根据表~\ref{tab:data2} 数据，以工作电流 $I_s$ 为横坐标、霍尔电压 $U_H$ 为纵坐标作图。
数据点 $(1.00, 13.75)$、$(3.00, 41.25)$、$(5.00, 68.75)$、$(7.00, 96.25)$、$(9.00, 123.75)$ 同样呈现良好的线性关系。

线性拟合斜率：
\begin{equation}
    k_2 = \frac{\Delta U_H}{\Delta I_s} = \frac{123.75 - 13.75}{9.00 - 1.00} = \frac{110.0}{8.00} = 13.75\;\text{mV/mA}
\end{equation}

由 $U_H = K_H B I_s$，当 $I_M = 600\;\text{mA}$ 时 $B = k I_M = 2.615 \times 10^{-2} \times 0.6 = 1.569 \times 10^{-2}\;\text{T}$，有：
\begin{equation}
    K_H = \frac{k_2}{B} = \frac{13.75 \times 10^{-3}\;\text{V/mA}}{1.569 \times 10^{-2}\;\text{T}} \approx 0.876\;\text{mV/(mA}\cdot\text{T)}
\end{equation}

与前述计算结果一致，验证了霍尔效应的线性规律。

\subsection{误差分析}

实验误差主要来源于以下几个方面：

\begin{enumerate}[label=(\arabic*), leftmargin=3em]
    \item \textbf{不等势电压 $U_0$：} 由表~\ref{tab:data1} 数据可估算 $U_0 \approx (U_1+U_2+U_3+U_4)/4 \approx 0.05\;\text{mV}$，通过对称测量法已基本消除；
    \item \textbf{温差电效应 $U_E$：} 在非大电流（$I_s \leq 10\;\text{mA}$）、非强磁场条件下，$U_H \gg U_E$，其影响可忽略不计；
    \item \textbf{电磁铁气隙中磁场的非均匀性：} 霍尔元件放置位置的微小偏差可能导致实际磁场与标称值存在差异；
    \item \textbf{测量仪表精度：} 电压表的分辨率和读数误差，以及工作电流、励磁电流的稳定性；
    \item \textbf{温度影响：} 半导体载流子浓度随温度变化，导致霍尔灵敏度略有漂移。
\end{enumerate}

\noindent 总体而言，$U_H$--$I_M$ 和 $U_H$--$I_s$ 两组数据均呈现良好的线性关系，线性相关系数 $R^2$ 接近 1，验证了霍尔效应公式 $U_H = K_H I_s B$ 的正确性。

\section{思考题：}

\begin{enumerate}[label=\arabic* 、, leftmargin=4em, labelsep=0pt]

    \item \textbf{什么是霍尔效应？简述其物理本质。}

    \textbf{答：} 霍尔效应是指当固体导体（或半导体）置于磁场中且有电流通过时，导体内的载流子受到洛伦兹力作用而发生偏转，在垂直于电流和磁场的方向上产生横向电压（霍尔电压）的现象。其物理本质是运动的带电粒子在磁场中受洛伦兹力 $\mathbf{F}_L = -e\mathbf{v} \times \mathbf{B}$ 作用，导致电荷在样品两侧积累，建立霍尔电场 $E_H$，当电场力与洛伦兹力平衡时，形成稳定的霍尔电压 $U_H$。

    \item \textbf{对称测量法（异号法）可以消除哪些副效应？其原理是什么？}

    \textbf{答：} 对称测量法通过改变工作电流 $I_s$ 和磁场 $B$ 的方向，进行四次测量（$+I_s,+B$；$+I_s,-B$；$-I_s,-B$；$-I_s,+B$），再按 $U_H = (U_1-U_2+U_3-U_4)/4$ 计算霍尔电压。该方法可以消除：
    \begin{itemize}
        \item 不等势电压 $U_0$（仅与 $I_s$ 方向有关）；
        \item 能斯特效应引起的附加电压 $U_N$（仅与 $B$ 方向有关）；
        \item 里吉-勒迪克效应引起的附加电压 $U_{RL}$（仅与 $B$ 方向有关）。
    \end{itemize}
    其原理是利用各副效应电压对 $I_s$ 和 $B$ 方向的不同依赖关系，通过正反向组合测量使其在代数运算中相互抵消。

    \item \textbf{若磁场方向与霍尔元件平面不严格垂直，对测量结果有何影响？}

    \textbf{答：} 若磁场 $\mathbf{B}$ 与霍尔元件平面法线方向夹角为 $\theta$，则只有垂直于元件平面的分量 $B\cos\theta$ 对霍尔效应有贡献，测得的霍尔电压 $U_H = K_H I_s B\cos\theta$ 将小于磁场完全垂直时的值。因此实验中应确保霍尔元件平面与磁场方向严格垂直，否则将引入系统误差，使测量值偏小。

\end{enumerate}
\begin{enumerate}[label=\arabic* 、, leftmargin=4em, labelsep=0pt]
    \item \textbf{什么是不等势电压？如何消除它对霍尔电压测量的影响？} \\
    答：不等势电压$U_0$是由于霍尔元件的1、2两电极不在理想的等势面上，当有工作电流$I_s$通过时产生的电压降。其符号只与$I_s$的方向有关，与磁场$\mathbf{B}$的方向无关。可通过改变$I_s$的方向（即对称测量法）予以消除。

    \item \textbf{对称测量法为什么能够消除多种副效应？有哪些副效应无法通过此方法消除？} \\
    答：对称测量法通过改变$I_s$和$B$的方向，进行四种组合的测量，利用$U_0$只与$I_s$有关、$U_N$和$U_{RL}$只与$B$有关的特点，通过代数运算$U_H=(U_1-U_2+U_3-U_4)/4$消除$U_0$、$U_N$、$U_{RL}$。但温差电效应$U_E$与$I_s$和$B$的关系跟$U_H$相同，不能通过对称测量法消除，不过其值很小，在非大电流、非强磁场下可忽略。

    \item \textbf{若霍尔元件的材料由N型变为P型半导体，霍尔电压的极性会如何变化？} \\
    答：N型半导体的载流子是电子（带负电），P型半导体的载流子是空穴（带正电）。在相同方向的电流和磁场下，正负载流子受到的洛伦兹力方向相同，但由于载流子电荷符号相反，积累的电荷极性也相反，因此霍尔电压的极性会发生反转。
\end{enumerate}
